\section*{TÓM TẮT ĐỒ ÁN}
\phantomsection\addcontentsline{toc}{section}{\numberline {}TÓM TẮT ĐỒ ÁN}

Trong hoạt động bán lẻ và kinh doanh dịch vụ hiện đại, khả năng thấu hiểu hành vi và đo lường trải nghiệm khách hàng (Customer Experience - CX) là yếu tố then chốt quyết định sự gắn kết của khách hàng và hiệu quả kinh doanh của doanh nghiệp. Các phương pháp khảo sát truyền thống (như bảng câu hỏi CSAT/NPS sau mua hàng) thường gặp hạn chế lớn do tính xâm lấn, độ trễ thông tin cao và tỷ lệ phản hồi thấp. Việc ứng dụng thị giác máy tính nhằm tự động ghi nhận và phân tích phản ứng cảm xúc của khách hàng tại các điểm chạm thực tế mang lại một hướng tiếp cận đột phá, cung cấp dữ liệu tức thời mà không làm gián đoạn hành trình mua sắm tự nhiên của khách hàng.

Đồ án này nghiên cứu và phát triển \textbf{``Hệ thống phân tích trải nghiệm và hành vi khách hàng dựa trên nhận dạng biểu cảm khuôn mặt tích hợp CRM''}. Hệ thống bao gồm hai năng lực cốt lõi được liên kết chặt chẽ:
\begin{enumerate}
    \item \textbf{Nhận dạng biểu cảm khuôn mặt thời gian thực (Facial Expression Recognition - FER):} Xây dựng pipeline xử lý hình ảnh bao gồm phát hiện, căn chỉnh khuôn mặt, đánh giá chất lượng ảnh đầu vào (độ mờ, độ sáng) và phân loại biểu cảm theo 7 lớp chuẩn quốc tế (\textit{Angry, Disgust, Fear, Happy, Sad, Surprise, Neutral}) với độ trễ thấp phục vụ môi trường thực tế.
    \item \textbf{Tổng hợp chuỗi thời gian và Phân tích trải nghiệm khách hàng (Customer Analytics):} Thiết kế giải thuật làm mịn theo chuỗi thời gian (Temporal Smoothing) và phân tích phiên tương tác ẩn danh theo từng dấu vết cục bộ (\textit{local track}) tại 6 khu vực không gian trong cửa hàng (Cổng vào, Khu vực chờ, Quầy tư vấn, Khu sản phẩm, Quầy thanh toán, Cổng ra). Hệ thống suy luận trạng thái trải nghiệm bán lẻ đa tín hiệu, trích xuất ma trận chuyển dịch cảm xúc (\textit{Expression Transitions}) và tính toán các chỉ số kinh doanh thông minh (BI) như Tỷ lệ biểu cảm tích cực/tiêu cực, Độ biến thiên cảm xúc ($\Delta$) trước và sau mua hàng, cùng hiệu suất tư vấn của nhân viên.
\end{enumerate}

Hệ thống được thiết kế theo kiến trúc Microservices phân tách độc lập giữa dịch vụ AI (FastAPI), phân hệ nghiệp vụ quản trị và phân tích dữ liệu (Java Spring Boot, PostgreSQL) và bảng điều khiển trực quan (Next.js Clean Architecture). Kết quả thực nghiệm khẳng định tính ổn định, khả năng đáp ứng thời gian thực và giá trị ứng dụng thực tiễn cao của hệ thống trong việc hỗ trợ doanh nghiệp tối ưu hóa quy trình dịch vụ và nâng cao chất lượng trải nghiệm khách hàng.

\newpage
\cleardoublepage